% mcm-paper source binding (non-rendering metadata)
% Repository HEAD: 66f49838fe3060c9c0286cb5aa87613ac20b9fc3
% Section role: model evaluation and future improvements; no new numerical claims introduced.

\section{模型评价与改进}\label{sec:evaluation}

\subsection{模型优点}\label{subsec:model-strengths}

本文以统一架次可行性核为基础，逐步将单向出海、联合取送和有限机队多日排班纳入同一建模框架。问题一通过物理路线执行次数与人员流变量分离，避免将人员配置固化在候选列中；问题二进一步利用半开区间动态容量刻画同站座位即时复用；问题三则在此基础上引入时间可行架次列和具体飞机指派，使模型能够自然承接时间窗、周转和单机冲突约束。三问之间具有清晰的递进关系，公共规则可以复用，新增约束也具有明确物理含义。

在求解层面，候选路线生成与整数主问题相分离，使大规模组合空间能够通过“启发式构池+精确整数优化”的方式处理。路线的容量、燃油、加油和机场闭环等条件在进入主问题前即完成可行性筛选，从而减少无效变量；求解后再将聚合人员流回填为具体人员和实际架次，提高了模型结果向可执行运输方案转化的能力。

多目标处理采用词典序优化，而非人为设定加权系数，使飞机使用时间、人员在途时间、燃油、架次数和座位利用率之间的优先级保持清晰。特别地，问题三将非临时任务基准、临时任务接纳和次级指标细化分成A--C三个阶段，使“完成刚性任务”和“利用剩余能力提高临时任务满足率”在逻辑上相互分离。其中阶段B在固定时间预算和当前有限时间可行架次列池下实现160/160全部接纳，并获得相对间隙为0的严格认证，也增强了阶段性结论的可信度。

此外，三问最终交付方案均通过独立校验程序复核。验证过程脱离求解器内部对象，重新检查人员完整性、容量、燃油及路线等硬约束，问题三还进一步检查时间窗、具体飞机冲突和30 min周转；同时配合候选池与时间列密度敏感性实验，使结果的可行性、稳定性和认证边界均有对应证据支撑。

\subsection{模型不足与改进}\label{subsec:model-limitations}

首先，问题一、问题二以及问题三的阶段A和阶段C均建立在有限候选路线或时间可行架次列池上，候选集由启发式规则生成，因此当前结果不能等同于完整可行空间上的全局最优。后续可引入列生成或定价思想，根据主问题对偶信息动态产生更有价值的路线，并进一步结合Branch-and-Price或自适应扩池策略，减少人工限定候选池所带来的最优性缺口。

其次，问题三采用临界起飞时刻离散化以控制模型规模。时间列密度实验表明，更密的时间采样能够提供更丰富的排班组合，但也显著增加整数搜索负担。后续可采用自适应时间离散方法，仅在紧张时间窗或高冲突区域增密候选时刻；也可将7日排班分解为日级或滚动时域子问题，再通过跨日资源协调逐步修正，以提高大规模时空调度的求解效率。

再次，当前词典序优化的后级目标依赖上一级获得的当前解，其中问题三阶段C尚未完成全部五级最优性认证，座位利用率阶段也未形成新的正式改进结果。后续可增加后级目标求解预算，并结合$\varepsilon$-约束法或Pareto前沿分析，在保持主目标允许微小退让的条件下考察人员时间、燃油和利用率之间的替代关系，从而获得更丰富的决策信息。

最后，现有敏感性分析主要覆盖问题一的加油候选保留数和问题三的时间列密度，问题二尚未针对路线池构造参数开展独立扰动实验。后续可系统改变路线池规模、成对插入候选范围和求解时间预算，并采用多随机种子重复求解，以进一步检验联合取送结果的稳定性。若进一步考虑实际运行中的天气、需求变化、飞行时间波动或临时延误，可在现有确定性模型上增加时间缓冲、情景约束或随机/鲁棒优化机制。已有研究分别采用情景化样本平均近似处理带时间窗取送问题中的随机需求\cite{ghilas2016stochastic}，以及采用两阶段随机规划刻画异构飞行器随机速度与时间窗约束\cite{jia2018}，可为后续不确定性扩展提供参考。
